% ═══════════════════════════════════════════════════════════════
%  FAS辅助主动RIS无人机安全通信系统模型 — TikZ 矢量图
%  编译方式: xelatex system_model_tikz.tex (支持中文)
%  输出: system_model_tikz.pdf
% ═══════════════════════════════════════════════════════════════
\documentclass[border=5pt]{standalone}
\usepackage{xeCJK}
\usepackage{tikz}
\usetikzlibrary{calc,arrows.meta,positioning,shapes.geometric,decorations.pathmorphing,shadows.blur,fit,backgrounds}

\setCJKmainfont{Microsoft YaHei}

% ── 颜色定义 ──
\definecolor{uavblue}{HTML}{2C5AA0}
\definecolor{uavlight}{HTML}{5C8BBF}
\definecolor{fasgreen}{HTML}{1A8A5C}
\definecolor{faslight}{HTML}{4CAF50}
\definecolor{risorange}{HTML}{D4760A}
\definecolor{rislight}{HTML}{FF9800}
\definecolor{userblue}{HTML}{2196F3}
\definecolor{userdark}{HTML}{1565C0}
\definecolor{evered}{HTML}{E53935}
\definecolor{evedark}{HTML}{B71C1C}
\definecolor{bldgray}{HTML}{CFD8DC}
\definecolor{blddark}{HTML}{90A4AE}
\definecolor{groundbrown}{HTML}{8D6E63}
\definecolor{skyblue}{HTML}{E3F2FD}

\begin{document}
\begin{tikzpicture}[
    >=Stealth,
    font=\sffamily\small,
    % ── 样式定义 ──
    uav/.style={circle, fill=uavblue, minimum size=1.2cm, draw=white, line width=1.5pt, blur shadow={shadow blur steps=5, shadow xshift=0.5pt, shadow yshift=-0.5pt}},
    rotor/.style={ellipse, fill=uavlight, minimum width=0.7cm, minimum height=0.25cm, draw=uavblue, line width=0.8pt, opacity=0.7},
    fas/.style={rounded corners=2pt, fill=fasgreen, minimum width=2.8cm, minimum height=0.3cm, draw=white, line width=1pt},
    port_active/.style={circle, fill=faslight, minimum size=0.18cm, draw=white, line width=0.8pt},
    port_inactive/.style={circle, fill=blddark!70, minimum size=0.14cm, draw=white, line width=0.5pt},
    ris/.style={rounded corners=3pt, fill=risorange, minimum width=2.4cm, minimum height=0.6cm, draw=white, line width=1.5pt},
    ris_cell/.style={rectangle, fill=rislight!70, draw=risorange, line width=0.3pt, minimum width=0.22cm, minimum height=0.15cm},
    user/.style={circle, fill=userblue, minimum size=0.8cm, draw=white, line width=1.2pt},
    eve/.style={circle, fill=evered, minimum size=0.8cm, draw=white, line width=1.2pt},
    bld/.style={rectangle, fill=bldgray, draw=blddark, line width=0.8pt},
    win/.style={rectangle, fill=userblue!30, draw=blddark!50, line width=0.3pt},
    sig/.style={->, line width=1.2pt, color=userblue},
    sig_dash/.style={->, line width=1.2pt, color=evered, dashed, dash pattern=on 4pt off 3pt},
    ch_label/.style={font=\footnotesize\itshape, fill=white, fill opacity=0.85, inner sep=1.5pt, rounded corners=1pt},
]

% ═══════════════════════════════════════════
%  1. 背景
% ═══════════════════════════════════════════

% 天空
\fill[skyblue, opacity=0.3] (-1,-0.5) rectangle (11,8);

% 地面
\fill[groundbrown, opacity=0.25] (-1,-0.5) rectangle (11,1.2);
\foreach \x in {-0.8,-0.4,...,10.8}{
    \draw[groundbrown, opacity=0.3, line width=0.3pt] (\x,0.3) -- (\x+0.3,0);
}

% ═══════════════════════════════════════════
%  2. 建筑物 (背景)
% ═══════════════════════════════════════════

% 建筑1 (左侧高楼)
\node[bld, minimum width=1cm, minimum height=3.5cm] (bld1) at (0.8,3.5) {};
\foreach \y in {1.6,2.1,...,6.0}{
    \node[win] at ($(bld1.north west)+(0.15,-\y+3.5)$) {};
    \node[win] at ($(bld1.north west)+(0.55,-\y+3.5)$) {};
    \node[win] at ($(bld1.north west)+(0.85,-\y+3.5)$) {};
}

% 建筑2 (中间矮楼)
\node[bld, minimum width=0.8cm, minimum height=2.5cm] (bld2) at (2.5,2.8) {};
\foreach \y in {0.6,1.1,...,3.6}{
    \node[win] at ($(bld2.north west)+(0.12,-\y+2.5)$) {};
    \node[win] at ($(bld2.north west)+(0.48,-\y+2.5)$) {};
}

% 建筑3 (右侧高楼)
\node[bld, minimum width=1.1cm, minimum height=3.8cm] (bld3) at (9.5,3.7) {};
\foreach \y in {1.8,2.3,...,6.8}{
    \node[win] at ($(bld3.north west)+(0.15,-\y+3.8)$) {};
    \node[win] at ($(bld3.north west)+(0.55,-\y+3.8)$) {};
    \node[win] at ($(bld3.north west)+(0.95,-\y+3.8)$) {};
}

% ═══════════════════════════════════════════
%  3. 无人机 (UAV)
% ═══════════════════════════════════════════

% 旋翼
\node[rotor] (r1) at (5.2,7.3) {};
\node[rotor] (r2) at (6.6,7.3) {};
\node[rotor] (r3) at (5.2,6.3) {};
\node[rotor] (r4) at (6.6,6.3) {};

% 旋翼臂
\draw[uavblue, line width=1.5pt, opacity=0.7] (r1) -- (5.9,6.8);
\draw[uavblue, line width=1.5pt, opacity=0.7] (r2) -- (5.9,6.8);
\draw[uavblue, line width=1.5pt, opacity=0.7] (r3) -- (5.9,6.8);
\draw[uavblue, line width=1.5pt, opacity=0.7] (r4) -- (5.9,6.8);

% 机身
\node[uav] (uav) at (5.9,6.8) {};
\node[white, font=\sffamily\bfseries\small] at (uav) {UAV};

% ═══════════════════════════════════════════
%  4. FAS 天线阵列
% ═══════════════════════════════════════════

% FAS 基板
\node[fas] (fas) at (5.9,6.1) {};

% FAS 端口
\node[port_active] (p3) at ($(fas.west)+(0.25,0)$) {};
\node[port_active] (p7) at ($(fas.west)+(1.25,0)$) {};
\node[port_active] (p10) at ($(fas.west)+(1.85,0)$) {};
\node[port_inactive] at ($(fas.west)+(0.05,0)$) {};
\node[port_inactive] at ($(fas.west)+(0.45,0)$) {};
\node[port_inactive] at ($(fas.west)+(0.65,0)$) {};
\node[port_inactive] at ($(fas.west)+(0.85,0)$) {};
\node[port_inactive] at ($(fas.west)+(1.05,0)$) {};
\node[port_inactive] at ($(fas.west)+(1.45,0)$) {};
\node[port_inactive] at ($(fas.west)+(1.65,0)$) {};
\node[port_inactive] at ($(fas.west)+(2.05,0)$) {};
\node[port_inactive] at ($(fas.west)+(2.55,0)$) {};

% FAS 标签
\node[font=\sffamily\tiny\bfseries, fasgreen, above=1pt of fas] {FAS ($N$=12端口)};

% ═══════════════════════════════════════════
%  5. 主动 RIS (建筑物上)
% ═══════════════════════════════════════════

\node[ris] (ris) at (5.9,3.0) {};

% RIS 单元格
\foreach \i in {0,...,7}{
    \node[ris_cell] at ($(ris.north west)+(0.2+\i*0.28,-0.12)$) {};
    \node[ris_cell] at ($(ris.north west)+(0.2+\i*0.28,-0.32)$) {};
}

% RIS 标签
\node[font=\sffamily\tiny\bfseries, risorange, above=2pt of ris] {主动 RIS ($M$=64 单元)};
\node[font=\sffamily\tiny, risorange!80, below=1pt of ris] {信号反射 + 人工噪声};

% ═══════════════════════════════════════════
%  6. 合法用户 (手机图标)
% ═══════════════════════════════════════════

% 用户1
\node[user] (u1) at (1.8,1.8) {};
\node[white, font=\sffamily\bfseries\footnotesize] at (u1) {U1};
\node[font=\sffamily\tiny\bfseries, userdark, below=3pt of u1] {User 1};

% 用户2
\node[user] (u2) at (10.0,1.8) {};
\node[white, font=\sffamily\bfseries\footnotesize] at (u2) {U2};
\node[font=\sffamily\tiny\bfseries, userdark, below=3pt of u2] {User 2};

% ═══════════════════════════════════════════
%  7. 窃听者
% ═══════════════════════════════════════════

\node[eve] (eve) at (5.9,1.8) {};
\node[white, font=\sffamily\bfseries\footnotesize] at (eve) {E};
\node[font=\sffamily\tiny\bfseries, evedark, below=3pt of eve] {Eavesdropper $p$};

% ═══════════════════════════════════════════
%  8. 信号路径
% ═══════════════════════════════════════════

% ── 直射链路: FAS → U1 ──
\draw[sig] (fas.south west) to[out=-120, in=100] node[ch_label, pos=0.4] {$\mathbf{h}^{U}_{U,k}$} (u1.north);

% ── 直射链路: FAS → U2 ──
\draw[sig] (fas.south east) to[out=-60, in=80] node[ch_label, pos=0.4] {$\mathbf{h}^{U}_{U,k}$} (u2.north);

% ── 直射链路: FAS → E (窃听, 红色虚线) ──
\draw[sig_dash] (fas.south) to[out=-100, in=90] node[ch_label, pos=0.4, text=evered] {$\mathbf{h}^{U}_{U,p}$} (eve.north);

% ── FAS → RIS (上行) ──
\draw[sig] (fas.south) to[out=-80, in=80] node[ch_label, pos=0.5] {$\mathbf{h}^{U}_{U,R}$} (ris.north);

% ── RIS → U1 (反射, 相长干涉) ──
\draw[sig] (ris.west) to[out=180, in=60] node[ch_label, pos=0.4] {$\mathbf{h}^{U}_{R,k}$} (u1.north east);

% ── RIS → U2 (反射, 相长干涉) ──
\draw[sig] (ris.east) to[out=0, in=120] node[ch_label, pos=0.4] {$\mathbf{h}^{U}_{R,k}$} (u2.north west);

% ── RIS → E (干扰, 人工噪声, 红色虚线) ──
\draw[sig_dash] (ris.south) to[out=-80, in=100] node[ch_label, pos=0.5, text=evered] {$\mathbf{h}^{U}_{R,p}$} (eve.north);

% ═══════════════════════════════════════════
%  9. 高度标注
% ═══════════════════════════════════════════

\draw[|<->|, line width=0.6pt, gray!70] (-0.3,1.5) -- node[font=\tiny, rotate=90, fill=white, fill opacity=0.8, inner sep=1pt] {$H=50$m} (-0.3,7.0);

% ═══════════════════════════════════════════
%  10. 图例
% ═══════════════════════════════════════════

\begin{scope}[shift={(8.8,0.5)}]
    \node[rounded corners=2pt, fill=white, draw=blddark!50, line width=0.5pt, minimum width=1.6cm, minimum height=1.2cm, opacity=0.95] at (0.6,0.5) {};
    \node[font=\sffamily\tiny\bfseries, above right] at (-0.1,1.0) {图例};

    % LOS
    \draw[sig, line width=1pt] (-0.05,0.75) -- (0.45,0.75);
    \node[font=\sffamily\tiny, right] at (0.55,0.75) {LOS};

    % virtual LOS
    \draw[sig, line width=1pt, dashed, dash pattern=on 3pt off 2pt] (-0.05,0.5) -- (0.45,0.5);
    \node[font=\sffamily\tiny, right] at (0.55,0.5) {virtual LOS};

    % Eve link
    \draw[sig_dash, line width=1pt] (-0.05,0.25) -- (0.45,0.25);
    \node[font=\sffamily\tiny, right] at (0.55,0.25) {Eavesdropper};

    % Active port
    \node[port_active, minimum size=0.12cm] at (0.2,0.0) {};
    \node[font=\sffamily\tiny, right] at (0.55,0.0) {FAS Active Port};
\end{scope}

% ═══════════════════════════════════════════
%  11. 标题
% ═══════════════════════════════════════════

\node[font=\sffamily\large\bfseries, above=0.3cm of uav, yshift=0.8cm] {FAS辅助主动RIS无人机安全通信系统模型};

% ═══════════════════════════════════════════
%  12. 区域标注
% ═══════════════════════════════════════════

% 信号增强区
\draw[userblue, dashed, line width=0.5pt, rounded corners=3pt, opacity=0.4]
    ($(u1.north west)+(-0.3,0.3)$) rectangle ($(u1.south east)+(0.3,-0.2)$);
\node[font=\tiny, userblue, opacity=0.6] at ($(u1.south)+(0,-0.35)$) {信号增强区};

% 干扰抑制区
\draw[evered, dashed, line width=0.5pt, rounded corners=3pt, opacity=0.4]
    ($(eve.north west)+(-0.3,0.3)$) rectangle ($(eve.south east)+(0.3,-0.2)$);
\node[font=\tiny, evered, opacity=0.6] at ($(eve.south)+(0,-0.35)$) {干扰抑制区};

\end{tikzpicture}
\end{document}
